\section{Discussion}
\label{sec:discussion}

\paragraph{The scope the evidence supports.}
The claim we support is bounded. For excerpts of IoT-style operational records, queried by a
question that names an entity whose value must not reach the model, routing each unit on
(sensitivity, relevance) and aliasing the sensitive-and-relevant cell preserves the task at
half the tokens and at zero measured leakage. The three variants in
Section~\ref{sec:analysis} delimit it: the rule is unnecessary when the question is topical,
redundant when anchor and answer coincide in a single record, and dependent on the refill
order when the answer lies outside the fallback window. We do not claim that it improves
ordinary compression on ordinary prompts, and the topical variant shows that it does not.

\paragraph{What the privacy claim covers.}
Leakage is measured as recoverability of the protected values from the exact prompt text,
using the pattern set of the published protocol (full value plus fragments). That is the
relevant quantity for a pipeline whose adversary is the model provider and whose prompts may
be retained. It is not a formal guarantee. We derive no differential-privacy
bound~\cite{dp} and no $k$-anonymity guarantee~\cite{kanon}; the alias map is assumed private,
and an adversary who obtains it recovers the values. We fix the alias assignment from a keyed
permutation precisely so that an adversary holding only the value set cannot invert the map
by rank, which is a real attack on unkeyed assignment, but that mitigation is narrower than a
guarantee. We also do not evaluate inference attacks---membership, attribute, or the
reconstruction of a value from its alias's co-occurrence pattern---and none of them is
addressed by the present evidence.

\paragraph{The refill order is a regularity.}
The measured benefit of keeping repetitive units follows from an anti-correlation between
redundancy and answer-bearing in this data. It is not derived from the definition of the
score, and a corpus whose answers are unique rather than repeated would not exhibit it. The
principled fix is to make the relevant region cover the answer instead of relying on the
ordering. In the two-hop variant the answer lies at a fixed distance of six sentences from
the anchor, so a window of that width would protect it explicitly and remove the dependence.
We have not run that configuration, and the present results should be read as depending on
the fallback width.

\paragraph{Cost.}
Prompt construction is lexical and model-free: one pass over the units to label them, one
pass to score redundancy (quadratic in the number of units in the worst case), and one pass
to fit the budget. It adds no model call. On the longest excerpt in the extension set
($19{,}838$ tokens uncompressed) construction completes in under a second on CPU. We did not
measure end-to-end wall-clock or memory of the reader, so the efficiency we claim concerns
prompt length rather than latency.

\paragraph{Threats to validity.}
The protected values are synthetic canaries injected into public corpora, because those
corpora cannot carry real personal data. Canary leakage is measurable exactly, but the
measured rate is a property of this construction. Both benchmark sizes are generated by our
own scripts from shared templates, so the questions are entity-resolving by construction and
the two-hop distance is fixed rather than sampled. The readers are 7B-parameter
instruction-tuned models, and we have no evidence about larger or API-served models. Finally,
the two benchmark sizes reuse case identifiers for different content, so a cross-set join by
identifier would silently mix records; we guard against this by verifying the uncompressed
arm against a per-set reference value, but the hazard is worth stating for anyone reusing the
artifacts.
